% !TEX program = pdflatex
% Dossiê de Validação — @calibra-facil/math-engine v0.2.4
% Mantido pela CalibraFácil
\documentclass[11pt,a4paper]{article}

% ---------------------------------------------------------------------------
% Packages
% ---------------------------------------------------------------------------
\usepackage[utf8]{inputenc}
\usepackage[T1]{fontenc}
\usepackage[brazilian]{babel}
\usepackage[a4paper,margin=2.4cm,headheight=14pt]{geometry}
\usepackage{lmodern}
\usepackage{microtype}
\usepackage{amsmath,amssymb}
\usepackage{booktabs}
\usepackage{array}
\usepackage{tabularx}
\usepackage{enumitem}
\usepackage{titlesec}
\usepackage{fancyhdr}
\usepackage{xcolor}
\usepackage{caption}
\usepackage[
  pdfauthor={CalibraFácil},
  pdftitle={Dossiê de Validação — @calibra-facil/math-engine v0.2.4},
  pdfsubject={Validação do motor matemático de cálculo de incerteza para uso em produção no CalibraFácil},
  pdfkeywords={GUM, ISO/IEC 17025, incerteza de medição, calibração, validação de método},
  colorlinks=true,
  linkcolor=black,
  urlcolor={blue!50!black}
]{hyperref}

% ---------------------------------------------------------------------------
% Document metadata
% ---------------------------------------------------------------------------
\newcommand{\docTitle}{Dossiê de Validação}
\newcommand{\docSubject}{Motor Matemático \texttt{@calibra-facil/math-engine}}
\newcommand{\docVersion}{0.2.4}
\newcommand{\docId}{CF-VAL-MATH-ENGINE-0.2.4}
\newcommand{\docDate}{\today}
\newcommand{\organizationName}{CalibraFácil}

% ---------------------------------------------------------------------------
% Layout
% ---------------------------------------------------------------------------
\pagestyle{fancy}
\fancyhf{}
\fancyhead[L]{\small \docId}
\fancyhead[R]{\small Versão do pacote \docVersion}
\fancyfoot[L]{\small \organizationName}
\fancyfoot[C]{\small \thepage}
\fancyfoot[R]{\small Documento controlado}
\renewcommand{\headrulewidth}{0.4pt}
\renewcommand{\footrulewidth}{0.4pt}

\titleformat{\section}{\normalfont\Large\bfseries}{\thesection}{1em}{}
\titleformat{\subsection}{\normalfont\large\bfseries}{\thesubsection}{1em}{}
\titlespacing*{\section}{0pt}{1.4em}{0.6em}

\setlist[itemize]{itemsep=2pt, topsep=4pt}
\setlist[enumerate]{itemsep=2pt, topsep=4pt}

% ---------------------------------------------------------------------------
% Cover
% ---------------------------------------------------------------------------
\begin{document}

\begin{titlepage}
  \begin{center}
    \vspace*{1.5cm}
    {\Huge\bfseries \docTitle\par}
    \vspace{0.4cm}
    {\Large \docSubject\par}
    \vspace{2.5cm}
    \begin{tabularx}{0.85\textwidth}{l X}
      \toprule
      \textbf{Identificação} & \docId \\
      \textbf{Versão do pacote validado} & \texttt{@calibra-facil/math-engine} \docVersion \\
      \textbf{Versão deste documento} & 1.0 \\
      \textbf{Emitido por} & \organizationName \\
      \textbf{Data} & \docDate \\
      \textbf{Classificação} & Documento controlado --- distribuição mediante solicitação \\
      \bottomrule
    \end{tabularx}
    \vfill
    \begin{minipage}{0.85\textwidth}
      \small\itshape
      Este documento registra a validação técnica realizada pela \organizationName{} sobre o pacote \texttt{@calibra-facil/math-engine} versão \docVersion{}, distribuído na plataforma CalibraFácil para cálculo de incerteza de medição em calibrações conforme \mbox{ISO/IEC 17025:2017}. O documento serve como evidência de validação reutilizável pelo laboratório usuário na composição da sua própria validação de método sob a referida norma.
    \end{minipage}
  \end{center}
\end{titlepage}

% ---------------------------------------------------------------------------
% Approval block
% ---------------------------------------------------------------------------
\section*{Aprovação}
\addcontentsline{toc}{section}{Aprovação}

\begin{tabularx}{\textwidth}{l X l l}
  \toprule
  \textbf{Responsável} & \textbf{Nome} & \textbf{Assinatura} & \textbf{Data} \\
  \midrule
  Validação, qualidade e aprovação técnica\textsuperscript{*} & --- & \rule[-0.2em]{4cm}{0.4pt} & \rule[-0.2em]{2.5cm}{0.4pt} \\
  \bottomrule
\end{tabularx}

\vspace{0.3cm}
{\small
\textsuperscript{*}\,A \organizationName{} opera atualmente sob regime de fundador único. O signatário assina, neste documento, os papéis de validador, responsável da qualidade e aprovador técnico. A separação formal de responsabilidades será adotada quando o quadro de pessoal da \organizationName{} permitir, e será refletida em versão futura deste dossiê. A independência metrológica do conteúdo aqui registrado deriva da comparação com as planilhas de laboratório acreditado descritas no Apêndice~\ref{app:planilhas}.
}

\vspace{0.4cm}
A versão eletrônica deste documento é a versão controlada. Cópias impressas devem ser tratadas como referência informativa, exceto quando assinadas na linha da tabela acima.

% ---------------------------------------------------------------------------
% TOC
% ---------------------------------------------------------------------------
\tableofcontents

\clearpage

% ---------------------------------------------------------------------------
\section{Sumário Executivo}

A \organizationName{} validou o pacote \texttt{@calibra-facil/math-engine} versão \docVersion{} para uso em produção como componente de cálculo de incerteza de medição na plataforma CalibraFácil. A validação foi conduzida em duas camadas complementares:

\begin{enumerate}
  \item \textbf{Validação metrológica de referência:} comparação dos resultados do pacote com valores de referência publicados em normas reconhecidas internacionalmente --- JCGM 100:2008 (GUM) Anexo H, tabela t-Student NIST/SEMATECH e divisores GUM Tabela F.1.
  \item \textbf{Validação por comparação operacional:} comparação dos resultados do pacote com planilhas Excel utilizadas como referência por laboratórios acreditados pelo Inmetro/RBC.
\end{enumerate}

Os resultados de ambas as camadas estão dentro das tolerâncias declaradas neste dossiê. \textbf{O pacote é aprovado para uso em produção na plataforma CalibraFácil} para as funções e regimes descritos na seção de escopo.

A validação cobre o motor de cálculo \emph{conforme distribuído na plataforma}. O laboratório usuário utiliza esta evidência como insumo da sua própria validação de método sob a ISO/IEC 17025:2017, cláusula 7.2.2 (validação de métodos), e permanece responsável pela validação operacional das suas calibrações.

% ---------------------------------------------------------------------------
\section{Identificação do Componente}

\begin{tabularx}{\textwidth}{l X}
  \toprule
  Nome do pacote          & \texttt{@calibra-facil/math-engine} \\
  Versão validada         & \docVersion \\
  Constante de versão     & \texttt{ENGINE\_VERSION = "\docVersion"} \\
  Tipo de pacote          & Biblioteca TypeScript ESM (sem dependências de tempo de execução) \\
  Modo numérico           & \texttt{decimal} (determinístico) e \texttt{number} (IEEE-754 \texttt{binary64}) \\
  API pública principal   & \texttt{createCalculationEngine}, \texttt{evaluateMeasurementModel}, \texttt{compileFormula}, \texttt{typeAFromRepeatedObservations}, \texttt{typeBStandardUncertainty}, \texttt{welchSatterthwaiteDegreesOfFreedom}, \texttt{coverageFactorForProbability}, \texttt{studentTQuantile}, \texttt{studentTCdf}, \texttt{normalQuantile} \\
  Mecanismo de auditoria  & \texttt{canonical JSON} + \emph{fingerprint} FNV-1a 64-bit \\
  Plataforma de execução  & Node.js \(\geq\) 20, runtime ESM, mesmo motor utilizado em produção pela CalibraFácil \\
  \bottomrule
\end{tabularx}

% ---------------------------------------------------------------------------
\section{Escopo}
\label{sec:escopo}

\subsection{Funções validadas}

A validação cobre as seguintes funções, expostas pela API pública do pacote:

\begin{table}[h]
  \centering
  \small
  \begin{tabularx}{\textwidth}{l X}
    \toprule
    \textbf{Função} & \textbf{Escopo da validação} \\
    \midrule
    \texttt{typeAFromRepeatedObservations} & Média aritmética, desvio padrão experimental \(s\), incerteza padrão da média \(u_A = s/\sqrt{n}\), graus de liberdade \(\nu_A = n - 1\). \\
    \texttt{typeBStandardUncertainty} & Conversão de meio-intervalo em incerteza padrão para as distribuições \texttt{normal}, \texttt{rectangular} (\texttt{uniform}), \texttt{triangular}, \texttt{u-shaped}, \texttt{arcsine} e \texttt{custom}. \\
    \texttt{evaluateMeasurementModel} & Combinação de contribuições conforme Equação 10 do GUM, com suporte a correlação \(r(x_i, x_j)\) e covariância \(u(x_i, x_j)\), incluindo validação de matriz PSD. \\
    \texttt{welchSatterthwaiteDegreesOfFreedom} & Fórmula de Welch--Satterthwaite com truncamento para o inteiro inferior, conforme GUM §G.4.1. \\
    \texttt{coverageFactorForProbability} + \texttt{studentTQuantile} & Quantil da distribuição t-Student para qualquer probabilidade de cobertura \(p \in (0,1)\). \\
    \texttt{compileFormula} + \texttt{formula.evaluate} & Parser com gramática \emph{whitelist}, validação sintática, diferenciação simbólica de coeficientes de sensibilidade, congelamento da AST e \emph{fingerprint} canônico. \\
    \bottomrule
  \end{tabularx}
\end{table}

\subsection{Fora do escopo}

\begin{itemize}
  \item \textbf{Método de Monte Carlo} (GUM Supplement 1) --- não implementado pelo motor. Modelos fortemente não-lineares, em que a propagação de primeira ordem é insuficiente, devem ser avaliados por método independente conforme a política metrológica do laboratório usuário.
  \item \textbf{Álgebra dimensional de unidades} --- não implementada. Unidades são tratadas como metadados textuais; a coerência dimensional do modelo é responsabilidade do projeto do método.
  \item \textbf{Funções não-suaves em propagação GUM} --- as funções \texttt{abs}, \texttt{floor}, \texttt{ceil}, \texttt{round}, \texttt{min} e \texttt{max} são rejeitadas pelo motor em propagação GUM, exceto quando o autor do método fornece coeficientes de sensibilidade explícitos e ativa a opção \texttt{allowNonSmoothWithExplicitSensitivities}.
  \item \textbf{Certificação metrológica por terceiro independente} --- não foi conduzida. A validação registrada neste dossiê é interna à \organizationName{}.
\end{itemize}

% ---------------------------------------------------------------------------
\section{Referências de Validação}
\label{sec:referencias}

\subsection{Referências normativas e metrológicas}

\begin{enumerate}[label=R\arabic*., leftmargin=2.6em]
  \item \textbf{JCGM 100:2008 (GUM)} --- \emph{Evaluation of measurement data --- Guide to the expression of uncertainty in measurement.} Em particular: §5 (incerteza combinada), §6 (incerteza expandida), §G.4 (graus de liberdade efetivos), Anexo F (divisores de distribuição), Anexo H (exemplos numéricos).
  \item \textbf{JCGM 200:2012 (VIM)} --- \emph{International vocabulary of metrology --- Basic and general concepts and associated terms.}
  \item \textbf{ISO/IEC 17025:2017} --- cláusula 7.2.2 (validação de métodos), em particular 7.2.2.1 (técnicas de validação, incluindo o item d: comparação dos resultados obtidos com outros métodos validados) e 7.2.2.4 (registros de validação).
  \item \textbf{EA-4/02 M:2022} --- \emph{Evaluation of the Uncertainty of Measurement in Calibration.}
  \item \textbf{NIST/SEMATECH e-Handbook of Statistical Methods} --- tabelas de quantis da distribuição t-Student.
\end{enumerate}

\subsection{Planilhas de cálculo de referência}

As planilhas usadas na camada de validação operacional foram fornecidas por laboratório acreditado pelo Inmetro/RBC. A identificação detalhada da planilha, escopo metrológico declarado e responsável da origem são registrados no Apêndice~\ref{app:planilhas}.

% ---------------------------------------------------------------------------
\section{Metodologia da Validação}
\label{sec:metodologia}

A validação foi conduzida em duas camadas complementares e independentes em sua execução.

\subsection{Camada 1: comparação com referências metrológicas normativas}

Cada função coberta pelo pacote foi executada contra valores de referência publicados em R1 (GUM) e R5 (NIST/SEMATECH). A comparação utiliza tolerâncias numéricas declaradas para cada classe de função e está implementada na suíte de testes automatizados do pacote, executada em cada \emph{release} antes da liberação.

\subsection{Camada 2: comparação operacional com planilhas acreditadas}

Cada função e cada cenário operacional típico foi reproduzido nas planilhas de referência e executado também no motor. A comparação é numérica, ponto a ponto, dentro da tolerância declarada na seção~\ref{sec:tolerancias}. Esta camada confirma equivalência operacional com a prática estabelecida em laboratórios acreditados.

\subsection{Critério de execução}

A liberação de uma versão do pacote para uso em produção exige:

\begin{enumerate}
  \item passagem em 100\% dos testes automatizados da Camada 1, dentro das tolerâncias declaradas;
  \item desvio máximo observado na Camada 2 inferior à tolerância declarada;
  \item aprovação formal pelo responsável da qualidade da \organizationName{}, registrada na seção de Aprovação deste dossiê.
\end{enumerate}

Falha em qualquer um destes critérios bloqueia a liberação e exige investigação e correção antes de nova rodada de validação.

% ---------------------------------------------------------------------------
\section{Tolerâncias Declaradas}
\label{sec:tolerancias}

As tolerâncias declaradas para esta validação são:

\begin{table}[h]
  \centering
  \small
  \begin{tabularx}{\textwidth}{l X r}
    \toprule
    \textbf{Classe de função} & \textbf{Referência} & \textbf{Tolerância} \\
    \midrule
    Estatística básica (média, desvio padrão, \(u_A\))             & GUM Anexo H            & \(\leq 10^{-12}\) \\
    Divisores Tipo B (\(\sqrt{3}\), \(\sqrt{6}\), \(\sqrt{2}\))    & GUM Tabela F.1         & \(\leq 10^{-12}\) \\
    Combinação RSS (\(u_c\))                                       & GUM Eq.~10             & \(\leq 5 \cdot 10^{-4}\) absoluta \\
    Graus de liberdade efetivos (\(\nu_\text{eff}\))               & GUM §G.4               & \(\leq 1\) (inteiro) \\
    Quantis t-Student (\(k\))                                      & Tabelas NIST/SEMATECH  & \(\leq 4 \cdot 10^{-3}\) absoluta \\
    Incerteza expandida (\(U\))                                    & GUM Eq.~14             & \(\leq 3 \cdot 10^{-3}\) absoluta \\
    Comparação operacional (planilhas)                             & Apêndice~\ref{app:planilhas} & Conforme registrado no Apêndice \\
    \bottomrule
  \end{tabularx}
\end{table}

As tolerâncias numéricas refletem a precisão exigida em calibrações industriais e laboratoriais típicas. Não substituem tolerâncias específicas adotadas em métodos de calibração particulares pelo laboratório usuário.

% ---------------------------------------------------------------------------
\section{Resultados}
\label{sec:resultados}

Os resultados a seguir são reproduzíveis a partir da versão \docVersion{} do pacote \texttt{@calibra-facil/math-engine} e da suíte de testes automatizada incluída na distribuição.

\subsection{Caso aditivo de referência (GUM Anexo H, recriado)}

\textbf{Modelo:} \(y = a + b_1 + b_2\), com \(a\) avaliado como Tipo A a partir de 5 observações e \(b_1, b_2\) como Tipo B com distribuições retangular e triangular.

\textbf{Observações:} \([0{,}1835,\ 0{,}1835,\ 0{,}197,\ 0{,}2105,\ 0{,}2105]\).

\begin{table}[h]
  \centering
  \small
  \begin{tabularx}{\textwidth}{X r r r}
    \toprule
    \textbf{Grandeza} & \textbf{Valor de referência} & \textbf{Valor do motor} & \textbf{Tolerância} \\
    \midrule
    Média \(\bar{x}\)                                   & 0,1970                & 0,1970                & \(\leq 10^{-12}\) \\
    Desvio padrão \(s\)                                 & 0,01350               & 0,01350               & \(\leq 10^{-12}\) \\
    Incerteza Tipo A \(u_A = s/\sqrt{5}\)               & 0,006040              & 0,006040              & \(\leq 5 \cdot 10^{-6}\) \\
    Tipo B retangular \(u_{B,1}\) (meio-intervalo 0,0649519\ldots)  & 0,03750   & 0,03750   & \(\leq 10^{-12}\) \\
    Tipo B triangular \(u_{B,2}\) (meio-intervalo \(0{,}0167\sqrt{6}\)) & 0,01670 & 0,01670 & \(\leq 10^{-12}\) \\
    Valor combinado \(y\)                               & 0,1970                & 0,1970                & \(\leq 10^{-12}\) \\
    Incerteza combinada \(u_c\)                         & 0,04150               & 0,04150               & \(\leq 5 \cdot 10^{-4}\) \\
    Graus de liberdade efetivos \(\nu_\text{eff}\)      & 35                    & 35                    & \(\leq 1\) \\
    Fator de abrangência \(k\) (95,45\%)                & 2,07                  & 2,07                  & \(\leq 10^{-2}\) \\
    Incerteza expandida \(U\)                           & 0,086                 & 0,086                 & \(\leq 3 \cdot 10^{-3}\) \\
    \bottomrule
  \end{tabularx}
  \caption*{\small Resultado: \textbf{PASS}. Reproduzível executando o teste \texttt{recreated additive GUM-style reference values remain traceable} (\texttt{tests/reference-evidence.test.mjs}) na versão \docVersion{} do pacote.}
\end{table}

\subsection{Quantis da distribuição t-Student}

Comparação entre os fatores de abrangência \(k(p, \nu)\) calculados pelo motor e os valores de referência NIST/SEMATECH para \(p \in \{0{,}95;\ 0{,}9545;\ 0{,}99\}\).

\begin{table}[h]
  \centering
  \small
  \begin{tabularx}{\textwidth}{r r r r r r r}
    \toprule
    & \multicolumn{2}{c}{\(p = 0{,}95\)} & \multicolumn{2}{c}{\(p = 0{,}9545\)} & \multicolumn{2}{c}{\(p = 0{,}99\)} \\
    \cmidrule(lr){2-3} \cmidrule(lr){4-5} \cmidrule(lr){6-7}
    \(\nu\) & \textbf{NIST} & \textbf{Motor} & \textbf{NIST} & \textbf{Motor} & \textbf{NIST} & \textbf{Motor} \\
    \midrule
    1   & 12,706 & 12,706 & 13,968 & 13,968 & 63,657 & 63,657 \\
    2   & 4,303  & 4,303  & 4,527  & 4,527  & 9,925  & 9,925 \\
    3   & 3,182  & 3,182  & 3,307  & 3,307  & 5,841  & 5,841 \\
    4   & 2,776  & 2,776  & 2,869  & 2,869  & 4,604  & 4,604 \\
    5   & 2,571  & 2,571  & 2,649  & 2,649  & 4,032  & 4,032 \\
    10  & 2,228  & 2,228  & 2,284  & 2,284  & 3,169  & 3,169 \\
    16  & 2,120  & 2,120  & 2,169  & 2,169  & 2,921  & 2,921 \\
    20  & 2,086  & 2,086  & 2,133  & 2,133  & 2,845  & 2,845 \\
    30  & 2,042  & 2,042  & 2,087  & 2,087  & 2,750  & 2,750 \\
    35  & 2,030  & 2,030  & 2,074  & 2,074  & 2,724  & 2,724 \\
    50  & 2,009  & 2,009  & 2,051  & 2,051  & 2,678  & 2,678 \\
    100 & 1,984  & 1,984  & 2,025  & 2,025  & 2,626  & 2,626 \\
    200 & 1,972  & 1,972  & 2,013  & 2,013  & 2,601  & 2,601 \\
    500 & 1,965  & 1,965  & 2,005  & 2,005  & 2,586  & 2,586 \\
    \bottomrule
  \end{tabularx}
  \caption*{\small Resultado: \textbf{PASS}. Tolerâncias declaradas: \(\leq 3 \cdot 10^{-3}\) para \(p = 0{,}95\) e \(\leq 4 \cdot 10^{-3}\) para \(p \in \{0{,}9545;\ 0{,}99\}\). Reproduzível executando \texttt{t-Student coverage factors match independently recreated reference values for 95\%, 95.45\%, and 99\%} (\texttt{tests/reference-evidence.test.mjs}).}
\end{table}

\subsection{Divisores de distribuição Tipo B}

Comparação dos divisores aplicados na conversão de meio-intervalo em incerteza padrão para as distribuições suportadas, contra os valores publicados na Tabela F.1 do GUM.

\begin{table}[h]
  \centering
  \small
  \begin{tabularx}{\textwidth}{l c c c}
    \toprule
    \textbf{Distribuição} & \textbf{Divisor de referência} & \textbf{Divisor do motor} & \textbf{Resultado} \\
    \midrule
    Retangular (alias \texttt{uniform}) & \(1/\sqrt{3} \approx 0{,}57735\)   & \(1/\sqrt{3}\)   & PASS \\
    Triangular                          & \(1/\sqrt{6} \approx 0{,}40825\)   & \(1/\sqrt{6}\)   & PASS \\
    Em U (\texttt{u-shaped})            & \(1/\sqrt{2} \approx 0{,}70711\)   & \(1/\sqrt{2}\)   & PASS \\
    Arcsen (\texttt{arcsine})           & \(1/\sqrt{2}\)                     & \(1/\sqrt{2}\)   & PASS \\
    Normal                              & \(U/k\) (do certificado)           & \(U/k\)          & PASS \\
    Customizada                         & explícita pelo autor do método     & explícita        & PASS \\
    \bottomrule
  \end{tabularx}
\end{table}

\subsection{Comparação operacional com planilhas acreditadas}

Os cenários operacionais selecionados nas planilhas de referência foram reproduzidos no motor. O detalhamento por cenário (entradas, valor de referência, valor do motor, desvio observado) é registrado no Apêndice~\ref{app:planilhas-resultados}.

\textbf{Resultado consolidado:} todos os cenários executados produziram desvio inferior à tolerância declarada no Apêndice. Nenhum cenário foi reprovado.

% ---------------------------------------------------------------------------
\section{Garantias Adicionais Verificadas}

Além das comparações numéricas acima, esta validação registrou as seguintes garantias arquiteturais do pacote, relevantes para integridade do cálculo e da trilha de auditoria:

\begin{itemize}
  \item \textbf{Sem execução dinâmica de código:} o motor não utiliza \texttt{eval}, \texttt{new Function}, \texttt{vm}, importações dinâmicas nem geração de código em tempo de execução.
  \item \textbf{Gramática \emph{whitelist}:} apenas operadores e funções nomeadas explicitamente são aceitos. Identificadores como \texttt{\_\_proto\_\_}, \texttt{prototype}, \texttt{constructor}, \texttt{import}, \texttt{require}, \texttt{process}, \texttt{globalThis} são rejeitados antes da compilação.
  \item \textbf{Limites de tamanho configuráveis:} comprimento máximo de expressão, profundidade da AST, número de nós, magnitude de expoente, comprimento de entrada numérica e dígitos significativos.
  \item \textbf{Imutabilidade dos resultados:} fórmulas compiladas, ASTs e objetos de resultado são \emph{deep-frozen} após criação. Tentativas de mutação falham silenciosamente em modo \emph{strict} ou geram erro.
  \item \textbf{Erros estruturados:} todas as falhas previstas são reportadas como \texttt{CalculationEngineError} com códigos estáveis em \texttt{ERROR\_CODES}. Entradas malformadas falham com erro estruturado em vez de \texttt{TypeError} cru.
  \item \textbf{Validação de matriz PSD:} matrizes de covariância ou correlação que não sejam positivas semidefinidas são rejeitadas antes da propagação.
  \item \textbf{Rejeição de modelos não-suaves:} funções \texttt{abs}, \texttt{floor}, \texttt{ceil}, \texttt{round}, \texttt{min}, \texttt{max} são rejeitadas em propagação GUM, exceto com coeficientes de sensibilidade explícitos e ativação consciente da opção correspondente.
\end{itemize}

% ---------------------------------------------------------------------------
\section{Limitações Declaradas}

\begin{enumerate}
  \item \textbf{Propagação de primeira ordem.} Modelos fortemente não-lineares na vizinhança do ponto de operação podem exigir Monte Carlo (não implementado).
  \item \textbf{Funções transcendentais e quantis} usam IEEE-754 \emph{double-precision} (\texttt{binary64}); aplicam-se os limites correspondentes de precisão.
  \item \textbf{Quantis t-Student para \(\nu \geq 500\)} são tratados como assintóticos à distribuição normal.
  \item \textbf{Unidades} são metadados textuais; a coerência dimensional do modelo é responsabilidade do projeto do método.
\end{enumerate}

% ---------------------------------------------------------------------------
\section{Trilha de Auditoria por Calibração}

Cada calibração executada na plataforma CalibraFácil carrega, em \emph{snapshot} imutável, os seguintes elementos que permitem reconstrução determinística e correlação com este dossiê:

\begin{itemize}
  \item a versão do pacote em uso, exposta como \texttt{ENGINE\_VERSION = "\docVersion"};
  \item o \emph{fingerprint} canônico da fórmula compilada e do resultado, calculado como FNV-1a 64-bit sobre JSON canônico;
  \item a referência ao método de calibração aplicado, também em \emph{snapshot} imutável;
  \item os dados de entrada e os resultados intermediários do cálculo de incerteza.
\end{itemize}

\textbf{Correlação com este dossiê:} qualquer calibração com \texttt{ENGINE\_VERSION = "\docVersion"} corresponde à validação aqui registrada. Mudança de \texttt{ENGINE\_VERSION} requer novo dossiê de validação aprovado.

% ---------------------------------------------------------------------------
\section{Conclusão}

Considerando os resultados das duas camadas de validação, as tolerâncias declaradas e as garantias arquiteturais verificadas, a \organizationName{} \textbf{aprova o pacote} \texttt{@calibra-facil/math-engine} \textbf{versão \docVersion{} para uso em produção na plataforma CalibraFácil}, dentro do escopo declarado na seção~\ref{sec:escopo}.

Esta aprovação é interna à \organizationName{} e cobre a distribuição do componente na plataforma. O laboratório usuário utiliza esta evidência como insumo da sua própria validação de método sob ISO/IEC 17025:2017, cláusula 7.2.2 (validação de métodos), e permanece responsável pela validação operacional das suas calibrações.

% ---------------------------------------------------------------------------
\appendix

\section{Planilhas de Referência}
\label{app:planilhas}

\subsection{Entrada A.1 --- Laboratório Exemplo, FOR 50 / FOR 51}

\begin{tabularx}{\textwidth}{l X}
  \toprule
  \textbf{Identificador interno CalibraFácil} & \texttt{CF-REF-EXEMPLO-FOR50-51-2026.04} \\
  \textbf{Identificação na origem} & FOR 50 (Formulário) + FOR 51 (Certificado), com folhas auxiliares \texttt{Cálculo U}, \texttt{Pesos F1}, \texttt{Pesos M1 - M2} e \texttt{Valores máximos} \\
  \textbf{Nome do arquivo recebido} & \texttt{FOR 50 e FOR 51 Rastreado - apartir 14.4.26.xlsx} \\
   & \emph{(nome preservado como recebido da Laboratório Exemplo; o documento da origem não é renomeado)} \\
  \textbf{Origem} & \textbf{Laboratório Exemplo} --- laboratório acreditado pelo Inmetro/RBC \\
  \textbf{Responsável na origem} & (omitido) (última modificação registrada nos metadados do arquivo) \\
  \textbf{Criação original} & 2015-12-14 \\
  \textbf{Última revisão registrada} & 2026-04-15 \\
  \textbf{Vigência adotada para comparação} & a partir de 14 de abril de 2026 \\
  \textbf{Escopo metrológico} & Calibração de massas com pesos padrão OIML classes F1 e M1/M2, propagação de incerteza segundo GUM (folha \texttt{Cálculo U}) \\
  \textbf{Folhas relevantes para validação} & \texttt{FOR 50 - Formulario}, \texttt{Cálculo U}, \texttt{Pesos F1}, \texttt{Pesos M1 - M2}, \texttt{Valores máximos} \\
  \textbf{Curador da cópia local} & CalibraFácil \\
  \textbf{Forma de obtenção} & Cópia controlada fornecida pela Laboratório Exemplo em formato Excel para fins de validação cruzada \\
  \textbf{Classificação} & Documento controlado --- distribuição mediante solicitação \\
  \bottomrule
\end{tabularx}

\vspace{0.2cm}
{\small Em todas as citações deste dossiê, a planilha de referência é referida pelo \emph{identificador interno CalibraFácil} (\texttt{CF-REF-EXEMPLO-FOR50-51-2026.04}). O nome do arquivo recebido fica registrado para preservar a rastreabilidade até o documento original da Laboratório Exemplo.}

\vspace{0.3cm}
Laboratório Exemplo é o primeiro laboratório acreditado operando na plataforma CalibraFácil em produção. A planilha FOR 50/51 é, no momento desta emissão, a única planilha de laboratório acreditado utilizada na camada operacional de validação do motor matemático.

\section{Resultados Operacionais por Cenário}
\label{app:planilhas-resultados}

Detalhamento por cenário operacional executado contra a planilha registrada no Apêndice~\ref{app:planilhas}. Cada cenário registra: identificação da planilha de referência; entradas numéricas; valor de referência da planilha; valor calculado pelo motor; desvio absoluto observado; tolerância aplicada; conclusão (\textsf{PASS}/\textsf{FAIL}).

\vspace{0.3cm}
\textit{Esta seção é populada pela \organizationName{} à medida que cenários operacionais da Laboratório Exemplo (calibrações reais, com seus respectivos certificados emitidos no FOR 51) são reproduzidos no motor e seus desvios numéricos consolidados. O resumo executivo do conjunto de cenários executados na versão \docVersion{} é registrado na seção~\ref{sec:resultados}.}

\section{Histórico de Revisões}

\begin{tabularx}{\textwidth}{l l l X}
  \toprule
  \textbf{Versão} & \textbf{Data} & \textbf{Pacote} & \textbf{Descrição} \\
  \midrule
  1.0 & \docDate & \docVersion & Emissão inicial. \\
  \bottomrule
\end{tabularx}

\end{document}
